% Identity Bedrock V14 — Pure minimal substrate, peer-review draft
% No PBAT, no physical presence extension.
% Compile: pdflatex -> pdflatex

\documentclass[11pt,a4paper]{article}

\usepackage[utf8]{inputenc}
\usepackage[T1]{fontenc}
\usepackage{lmodern}
\usepackage{microtype}
\usepackage{amsmath,amssymb,amsthm}
\usepackage{mathtools}
\usepackage{booktabs}
\usepackage{array}
\usepackage{enumitem}
\usepackage[colorlinks=true,linkcolor=black,citecolor=black,urlcolor=blue]{hyperref}
\usepackage{geometry}

\geometry{a4paper,margin=25mm}

\newtheorem{theorem}{Theorem}[section]
\newtheorem{definition}[theorem]{Definition}
\newtheorem{remark}[theorem]{Remark}
\newtheorem{corollary}[theorem]{Corollary}

\title{\textbf{Identity Bedrock: Operationalizing Rompel's Minimal\\
Cryptographic Substrate for Sovereign Identity}}
\author{Jean-Fran\c{c}ois Bouchard\\
  \small PrivateNess Network (co-founder, CTO)\\
  \small Emercoin (community contributor)\\
  \small \texttt{[contact via privateness.network]}}
\date{March 2026}

\begin{document}

\maketitle

\noindent\textbf{AI disclosure.} The author used AI language model assistance
for manuscript drafting and editing.  All technical contributions, system
design, implementation, theoretical framing, and binding insights described
herein originate with the author, drawing on nine years of operational
experience deploying Emercoin-based infrastructure.

\begin{abstract}
Identity Bedrock was not built from theory.  It was built from a requirement:
eliminate every institutional key beneath the individual.  Following that
requirement to its logical floor---asking what is the minimum cryptographic
primitive that supports public authentication---we arrived at one-way functions
as the irreducible foundation.  Rompel's 1990 proof that OWFs are necessary
and sufficient for existentially unforgeable signatures under adaptive
chosen-message attack~\cite{rompel1990} was found after that convergence.
He had already proved we were standing on bedrock.

The system separates identity cryptography from directory consensus via a
chain-agnostic, append-only binding layer, enabling portability across
heterogeneous substrates while preserving verifiability.  We instantiate this
construction using Emercoin Name--Value Storage (NVS) anchored to Bitcoin
via Auxiliary Proof-of-Work (AuxPoW), with a publicly verifiable deployment
operational since 2020.

We further show that additional system capabilities---naming, structured
record representation, content retrieval, deterministic execution, overlay
connectivity, and economic interaction---emerge as derived properties of this
identity substrate rather than as independent subsystems requiring their own
trust roots.

The contribution is not a new primitive but a composition of existing ones
into a system that minimises trust assumptions to the irreducible minimum
while remaining deployable on commodity infrastructure.
\end{abstract}

\noindent\textbf{Keywords:} one-way functions $\cdot$ digital signatures
$\cdot$ sovereign identity $\cdot$ blockchain identity $\cdot$
chain-agnostic architecture $\cdot$ append-only directory $\cdot$
post-quantum migration

%--------------------------------------------------------------------
\section{Introduction}
%--------------------------------------------------------------------

\subsection{The Problem: Institutional Capture}

Every deployed digital identity system places institutional keys beneath the
user: certificate authorities, identity providers, key registries, platform
operators.  These entities can revoke, deny, surveil, or control identity
without the user's consent.

The threat model motivating this work assumes an adversary with substantial
conventional resources, potential political backing, and full control of
traditional identity infrastructure.  The central question is: \emph{can
identity exist without institutional keys beneath the individual?}

\subsection{The Theoretical Answer: Rompel 1990}

Rompel proved that OWFs are both necessary and sufficient for
$\mathrm{SIG}_{\mathrm{EUF\text{-}CMA}}$ secure signature
schemes~\cite{rompel1990}.  This is the provable minimum---nothing strictly
weaker supports public authentication under standard assumptions.

Our thesis: fixing OWF as the cryptographic floor and adding only a
chain-agnostic append-only directory produces an identity construction that
minimises institutional capture to zero beneath the individual key.

\subsection{The Binding Insight}

This work follows a pattern established by prior binding innovations in
cryptographic systems: the value lies not in the invention of new primitives
but in recognising that existing primitives, bound in a specific way, produce
emergent properties none of them exhibits alone.  Each component of Identity
Bedrock---OWF-based signatures, append-only directories---exists independently
in prior literature.  The contribution is the binding and the formal argument
for the security properties that binding produces.

\subsection{Contributions}

\begin{enumerate}[label=\textbf{C\arabic*.},leftmargin=*]
  \item \textbf{Theory.}  Formalisation of a Rompel-grounded, chain-agnostic
    identity system with inline definitions and a formal authentication
    soundness theorem proved via explicit security reduction.

  \item \textbf{Systems.}  Architecture: OWF-based identity layer (Ed25519
    default; ML-DSA-65 migration target~\cite{fips204}) with pluggable
    directory substrate.  Implementation: Emercoin NVS anchored to Bitcoin
    via AuxPoW~\cite{emercoin} and NESS Fiber~\cite{onchain}.  Deployed and
    publicly verifiable since 2020.

  \item \textbf{Systemic emergence.}  Demonstration that naming, structured
    representation, content addressing, deterministic execution, overlay
    connectivity, and economic interaction all arise as derived properties of
    the identity substrate without additional trust roots.

  \item \textbf{Attack-cost analysis.}  Quantified cost of attacking both the
    cryptographic and directory components using the hack51
    framework~\cite{hack51}.
\end{enumerate}

\subsection{Paper Organisation}

Section~2 reviews background.  Section~3 defines the threat model.
Section~4 gives the formal model.  Section~5 presents the architecture.
Section~6 proves authentication soundness.  Section~7 describes derived
system properties.  Section~8 covers the concrete instantiation and
attack-cost analysis.  Section~9 demonstrates the deployed system.
Section~10 discusses limitations and future work.

%--------------------------------------------------------------------
\section{Background}
%--------------------------------------------------------------------

\subsection{One-Way Functions and Digital Signatures}

A function $f : \{0,1\}^* \to \{0,1\}^*$ is a \emph{one-way function} (OWF)
if it is computable in polynomial time and for every probabilistic
polynomial-time (PPT) adversary $\mathcal{A}$,
\[
  \Pr\bigl[f(x') = f(x) \;\big|\; x \gets \{0,1\}^\lambda,\;
    x' \gets \mathcal{A}(1^\lambda, f(x))\bigr] \leq \mathrm{negl}(\lambda).
\]

Rompel~\cite{rompel1990} proved, completing the programme of
Naor--Yung~\cite{naoryung1989} and later Haitner et
al.~\cite{haitner2010}, that OWFs are \emph{necessary and sufficient} for
EUF-CMA-secure digital signature schemes:
\[
  \mathrm{OWF} \iff \mathrm{SIG}_{\mathrm{EUF\text{-}CMA}}.
\]

This equivalence means that the existence of any EUF-CMA-secure scheme
already commits the system to exactly one hardness assumption---the existence
of OWFs---and that no weaker assumption is sufficient.  Identity Bedrock
operationalises this as a design principle: the individual's keypair is the
unique cryptographic root of identity.

\subsection{Traditional Identity Architectures}

Public-Key Infrastructure (PKI) / X.509~\cite{rfc5280} places a CA hierarchy
beneath the individual key, introducing revocation oracles and CA compromise
as additional failure modes.  Federated identity systems (OAuth 2.0, OpenID
Connect) delegate identity to a provider that can de-authenticate users
unilaterally.  Even in blockchain settings, governance multisigs and
administrative keys often sit below the individual in the trust stack.  These
architectural choices add assumptions beyond OWFs without increasing the
cryptographic security of authentication itself.

\subsection{Bitcoin AuxPoW and Merged Mining}

Auxiliary Proof-of-Work allows Emercoin to be merge-mined with Bitcoin,
inheriting a portion of Bitcoin's security for non-monetary applications.
Khovayko~\cite{hack51} introduces the \emph{hack51 coefficient}---defined as
$\text{attack\_cost} / \text{market\_cap}$---as a numerically measurable
metric for 51\% attack resistance.  Under this metric, Emercoin's triple
consensus (PoS + PoW + merged mining) produces a coefficient significantly
exceeding 1.0: attacking Emercoin requires resources approaching five times
its market capitalisation.  We apply the hack51 framework in
Section~\ref{sec:instantiation}.

%--------------------------------------------------------------------
\section{Threat Model}
%--------------------------------------------------------------------

\subsection{Adversary Capabilities}

\textbf{Assumed:} Substantial conventional resources and political backing;
full control of traditional identity infrastructure (CAs, IdPs, KYC
registries); ability to outlaw software, seize property, and coerce
individuals; classical (polynomial-time) computational adversary.

\textbf{Not assumed:} Cannot break OWF under standard complexity assumptions;
cannot mount a sustained 51\% attack on Bitcoin (quantified in
Section~\ref{sec:instantiation}); cannot simultaneously compromise distributed
infrastructure globally.

\textbf{Explicit scope.}  This work addresses classical adversaries.
Post-quantum migration to ML-DSA (FIPS~204~\cite{fips204}) is discussed in
Section~\ref{sec:discussion}.

\subsection{Design Requirements}

\begin{description}[leftmargin=*,style=nextline]
  \item[R1 --- Mathematical irreducibility.]  OWF as cryptographic floor.
    No strictly weaker primitive beneath.
  \item[R2 --- No institutional keys.]  Individual key is sole cryptographic
    root of trust.  No CA, IdP, or registry key beneath the user.
  \item[R3 --- Chain-agnostic.]  Identity cryptography independent of any
    specific directory consensus.  Portable across append-only substrates.
  \item[R4 --- Quantifiable security.]  Attack costs explicit and measurable.
  \item[R5 --- Operational verification.]  Deployed system, not theoretical
    proposal.  Public verification of all claims.
\end{description}

%--------------------------------------------------------------------
\section{Formal Model}
%--------------------------------------------------------------------

Let $\lambda \in \mathbb{N}$ denote the security parameter and
$\mathrm{negl}(\lambda)$ a negligible function.

\begin{definition}[Digital Signature Scheme]
  A digital signature scheme is a triple $\Sigma = (\mathsf{KeyGen},
  \mathsf{Sign}, \mathsf{Verify})$ of PPT algorithms where:
  \begin{itemize}
    \item $\mathsf{KeyGen}(1^\lambda) \to (sk, pk)$,
    \item $\mathsf{Sign}(sk, m) \to \sigma$ for $m \in \{0,1\}^*$,
    \item $\mathsf{Verify}(pk, m, \sigma) \to \{0,1\}$.
  \end{itemize}
  $\Sigma$ is \emph{correct} if $\mathsf{Verify}(pk, m, \mathsf{Sign}(sk,m))
  = 1$ for all $(sk,pk) \gets \mathsf{KeyGen}(1^\lambda)$ and all $m$.
\end{definition}

\begin{definition}[EUF-CMA Security]
  $\Sigma$ is \emph{existentially unforgeable under adaptive chosen-message
  attack} if for every PPT $\mathcal{A}$,
  \[
    \Pr\bigl[\mathsf{EUF\text{-}CMA}(\mathcal{A},\Sigma,\lambda) = 1\bigr]
    \leq \mathrm{negl}(\lambda),
  \]
  where $\mathsf{EUF\text{-}CMA}(\mathcal{A},\Sigma,\lambda)$ is the standard
  game: $(sk,pk) \gets \mathsf{KeyGen}(1^\lambda)$; $\mathcal{A}^{\mathsf{Sign}(sk,\cdot)}(pk)$
  outputs $(m^*, \sigma^*)$; $\mathcal{A}$ wins if $\mathsf{Verify}(pk,
  m^*, \sigma^*) = 1$ and $m^*$ was never queried to the signing oracle.
\end{definition}

\begin{definition}[Public Directory]
  A \emph{public directory} $\mathcal{D}$ is a set of pairs $(id, pk) \in
  \mathcal{ID} \times \mathcal{PK}$ with membership predicate
  \[
    \mathsf{Bind}(id, pk) = 1 \iff (id, pk) \in \mathcal{D}.
  \]
  $\mathcal{D}$ is \emph{append-only} if for every $(id, pk)$ once inserted,
  there is no efficient procedure that can cause $\mathsf{Bind}(id, pk)$ to
  become $0$.
\end{definition}

\begin{definition}[$\varepsilon_{\mathrm{dir}}$-Finality]
  An append-only directory $\mathcal{D}$ has \emph{$\varepsilon_{\mathrm{dir}}$-finality}
  if for any PPT adversary $\mathcal{A}$ and any entry $(id, pk)$ written to
  $\mathcal{D}$ at depth $d$,
  \[
    \Pr\bigl[\mathcal{A}\ \text{causes } \mathsf{Bind}(id,pk) = 0\bigr]
    \leq \varepsilon_{\mathrm{dir}}(\lambda, d).
  \]
  As confirmation depth $d \to \infty$, $\varepsilon_{\mathrm{dir}} \to 0$.
\end{definition}

\begin{definition}[Authentication Predicate]
  Given directory $\mathcal{D}$, the \emph{authentication predicate} is
  \[
    \mathsf{Auth}(id, c, \sigma) := 1 \iff
    \exists\, pk :\;
    \mathsf{Bind}(id, pk) = 1
    \;\wedge\;
    \mathsf{Verify}(pk, c, \sigma) = 1.
  \]
\end{definition}

%--------------------------------------------------------------------
\section{Architecture}
%--------------------------------------------------------------------

\subsection{Layer~1: Identity Cryptography}

A party generates $(sk, pk) \gets \mathsf{KeyGen}(1^\lambda)$ from a
high-entropy seed (OS entropy combined with an application-layer PRNG, e.g.,
UHEPRNG~\cite{uheprng}).  The public key $pk$, together with a chosen
identifier $id \in \mathcal{ID}$, constitutes the identity object
$\mathcal{I} = (id, pk)$.  Authentication to any verifier who can evaluate
$\mathsf{Bind}$ is then a standard challenge--response:

\begin{enumerate}[leftmargin=*]
  \item Verifier samples $c \gets \{0,1\}^\lambda$ and sends $c$ to the prover.
  \item Prover computes $\sigma \gets \mathsf{Sign}(sk, c)$ and returns $\sigma$.
  \item Verifier accepts iff $\mathsf{Auth}(id, c, \sigma) = 1$.
\end{enumerate}

No certificate, session token, or third-party oracle appears in this
exchange.

\subsection{Layer~2: Append-Only Directory}

The directory $\mathcal{D}$ binds identifiers to public keys.  The
requirements are exactly those of Definition~3: global verifiability,
persistence, and $\varepsilon_{\mathrm{dir}}$-finality.  The identity layer
and directory layer are decoupled: any substrate satisfying Definition~3 is a
valid instantiation.  Section~\ref{sec:instantiation} gives a concrete one.

%--------------------------------------------------------------------
\section{Security Analysis}
\label{sec:security}
%--------------------------------------------------------------------

\begin{theorem}[Authentication Soundness]
  \label{thm:soundness}
  Let $\Sigma$ be EUF-CMA secure with advantage bound $\varepsilon_{\mathrm{sig}}$,
  and let $\mathcal{D}$ be an append-only directory with
  $\varepsilon_{\mathrm{dir}}$-finality.  For any PPT adversary $\mathcal{A}$
  that does not hold $sk$, the probability that $\mathcal{A}$ causes a
  verifier to accept as identity $id$ is at most
  $\varepsilon_{\mathrm{sig}}(\lambda) + \varepsilon_{\mathrm{dir}}(\lambda, d)$.
\end{theorem}

\begin{proof}
Suppose adversary $\mathcal{A}$ wins the authentication game: it produces
valid $(id, c, \sigma)$ without holding $sk$.

\textbf{Reduction.}  Construct forger $\mathcal{B}$ against $\Sigma$.
$\mathcal{B}$ receives challenge public key $\hat{pk}$.  It simulates the
directory for $\mathcal{A}$ by answering $\mathsf{Bind}$ queries for $id$
with $\hat{pk}$.\footnote{$\mathcal{B}$ simulates the directory locally; it
does not write to any real ledger.}  We consider two mutually exclusive cases.

\textbf{Case 1: directory integrity holds} (probability $\geq 1 -
\varepsilon_{\mathrm{dir}}$).  The only $pk$ satisfying
$\mathsf{Bind}(id,pk)=1$ is $\hat{pk}$.  When $\mathcal{A}$ triggers
authentication, $\mathcal{B}$ forwards the EUF-CMA game challenge $c$ to
$\mathcal{A}$.  Since $c$ is chosen fresh after $\mathcal{A}$ commits its
strategy and was never queried to the legitimate $\mathsf{Sign}(sk,\cdot)$
oracle, any valid $\sigma$ returned by $\mathcal{A}$ is an existential
forgery.  By EUF-CMA security of $\Sigma$:
\[
  \Pr[\mathcal{A}\ \text{wins} \mid \text{Case 1}]
  \leq \varepsilon_{\mathrm{sig}}(\lambda).
\]

\textbf{Case 2: directory integrity fails} (probability $\leq
\varepsilon_{\mathrm{dir}}$).  $\mathcal{A}$ replaces $(id, \hat{pk})$ with
$(id, pk')$ for a key it controls.  This event occurs with probability at
most $\varepsilon_{\mathrm{dir}}(\lambda,d)$ by Definition~4.  Directory
attacks shift cost to the directory layer, quantified in
Section~\ref{sec:instantiation}.

By a union bound:
$\Pr[\mathcal{A}\ \text{wins}] \leq
\varepsilon_{\mathrm{sig}} + \varepsilon_{\mathrm{dir}}$,
both negligible under the stated assumptions.
\end{proof}

\begin{corollary}
  Any successful impersonation of identity $id$ requires either (i) breaking
  $\Sigma$ with advantage $> \varepsilon_{\mathrm{sig}}$, or (ii) corrupting
  the directory with cost exceeding the economic threshold implied by
  $\varepsilon_{\mathrm{dir}}$.  No institutional secret appears in either
  condition.
\end{corollary}

%--------------------------------------------------------------------
\section{Derived System Properties}
\label{sec:derived}
%--------------------------------------------------------------------

Once an identity substrate satisfying Definitions~1--5 exists, additional
system capabilities are obtainable as derived properties without introducing
new trust roots.

\subsection{Naming and Resolution}

Append-only bindings between human-readable names and public keys implement
decentralised naming.  Resolution is identity-check: to resolve name $n$,
retrieve $(n, pk)$ from $\mathcal{D}$ and verify $\mathsf{Bind}(n, pk) = 1$.
No hierarchical authority is required.  This subsumes the functionality of
DNS-over-blockchain proposals~\cite{handshake} without a separate trust
model.

\subsection{Structured Record Representation}

A canonical, append-only JSON schema---the WORM (Write-Once Read-Many)
format---defines typed records stored as directory values.  Deterministic
serialisation allows independent nodes to parse, index, and query records
without a shared schema authority.

\subsection{Content Addressing}

Content retrieval uses cryptographic hashes: a node stores a datum $m$ at
address $h = H(m)$ and any retriever verifies $H(m') = h$ before use.
Integrity is transport-independent~\cite{ipfs}.  Combined with the identity
layer, a publisher can sign a content address, binding retrieval integrity to
identity authenticity.

\subsection{Deterministic Execution}

A computation is modelled as a pure function $F$ applied to content-addressed
inputs: given inputs $\{h_i\}$ and a content-addressed program $h_F$, the
output $h_{\mathrm{out}} = H(F(m_1, \ldots, m_k))$ is deterministic and
independently reproducible by any node that can retrieve the inputs.
Verifiable execution reduces to checking that the published output hash
matches independent local evaluation.

\subsection{Overlay Connectivity}

Peer discovery and message routing can be implemented over identity-addressed
endpoints.  Because each endpoint is authenticated via $\mathsf{Auth}$, the
overlay inherits the security of the identity layer and does not require a
separate PKI for transport authentication.

\subsection{Economic Interaction}

Participation-based reward and payment models can be constructed using
identity-bound attribution.  Each economic action is signed by the actor's
$sk$, making contributions attributable and auditable without a centralised
ledger operator.

%--------------------------------------------------------------------
\section{Concrete Instantiation}
\label{sec:instantiation}
%--------------------------------------------------------------------

\subsection{Signature Scheme}

We instantiate $\Sigma$ with Ed25519~\cite{eddsa}, a Schnorr-variant
signature scheme over Curve25519 providing approximately $2^{128}$ security
against known attacks.  Key generation uses a multi-source high-entropy seed:
multiple 1024-byte blocks from \texttt{os.urandom} (1536 bits total internal
state), RC4OK output from Emercoin core RNG, CPU timing jitter, and hardware
identifiers, mixed via UHEPRNG~\cite{uheprng}.

\subsection{Directory: Emercoin NVS + AuxPoW}

The public directory $\mathcal{D}$ is implemented using Emercoin
Name--Value Storage (NVS)~\cite{emercoin}, a blockchain-native key--value
store.  Emercoin blocks are merge-mined against Bitcoin via Auxiliary
Proof-of-Work (AuxPoW)~\cite{hack51}, anchoring directory state to
Bitcoin's cumulative proof-of-work.

Identity records are written as WORM JSON objects under NVS keys of the form
\texttt{id:$\langle$identifier$\rangle$}.  Once written, records can be
renewed but not deleted, satisfying the append-only requirement of
Definition~3.

\subsection{Attack-Cost Analysis}

\begin{center}
\begin{tabular}{lll}
  \toprule
  \textbf{Attack vector} & \textbf{Requirement} & \textbf{Result} \\
  \midrule
  Break Ed25519    & Violate OWF assumption   & Infeasible ($2^{128}$ ops) \\
  Rewrite Emercoin NVS & PoS majority + AuxPoW 51\% & 5$\times$ EMC mkt.\ cap~\cite{hack51} \\
  Compromise NESS Fiber & Steal publisher keys & DoS only; no forgery \\
  Coerce enrolled user  & Physical / legal coercion & Out of cryptographic scope \\
  \bottomrule
\end{tabular}
\end{center}

\textbf{Signature forgery.}  Breaking Ed25519 requires solving the discrete
logarithm problem on Curve25519; no sub-exponential algorithm is known
($O(2^{128})$ operations, classically infeasible).

\textbf{Directory rewrite.}  Altering a finalised NVS entry requires a
$>50\%$ attack on Emercoin's AuxPoW chain.  Under the hack51
framework~\cite{hack51}, this costs resources approaching five times
Emercoin's market capitalisation.  Nuzzi et al.~\cite{nuzzi2024} place the
cost of attacking Bitcoin itself at USD 5--20B for a sustained attack,
providing a lower bound on the AuxPoW component alone.

\textbf{NESS Fiber key compromise.}  A stolen publisher key cannot rewrite
past blocks (append-only BoltDB makes history rewriting impossible).  The only
action available is a reset to genesis, which is immediately visible to all
participants.  Key compromise is thus maximally detectable.

%--------------------------------------------------------------------
\section{Practical Demonstration}
%--------------------------------------------------------------------

The system has been operational since 2020 on the PrivateNess Network.  The
canonical on-chain identity record~\cite{onchain}---genesis address, blockchain
public key, and genesis signature---is anchored in Emercoin NVS block~\#445669
(November~1, 2020, 22:16:45~UTC), transitively anchored to Bitcoin via AuxPoW,
and independently auditable at
\url{https://explorer.emercoin.com/dpo/PrivateNESS.Network} by any party
without trusting an intermediary.  This is not a paper describing what could
be built.  It is a paper describing what has been running since 2020.

\textbf{NESS Fiber.}  NESS Fiber uses Skycoin-type Fiber consensus~\cite{skycoin}
with announced publisher keys published in Emercoin NVS, creating a verifiable
trust chain: NESS Fiber security anchors to Emercoin NVS, which anchors to
Bitcoin via AuxPoW.  Any verifier can confirm the binding independently at the
Emercoin block explorer without relying on any intermediary.

\textbf{Resolution flow.}
\begin{enumerate}[leftmargin=*]
  \item Look up \texttt{id:$\langle$name$\rangle$} in Emercoin NVS;
    verify $\mathsf{Bind}(\mathit{name}, pk) = 1$.
  \item Retrieve associated content by hash; verify $H(m') = h$.
  \item Verify publisher signature: $\mathsf{Verify}(pk, h, \sigma) = 1$.
  \item Execute content-addressed program deterministically.
\end{enumerate}

Each step is independently verifiable by any node with access to the public
directory and standard hash functions.  No trusted intermediary participates.
The system operates on commodity hardware including RK3528-class single-board
computers and repurposed consumer devices, with a verified full-stack
deployment cost under 50~CAD~\cite{onchain}.

%--------------------------------------------------------------------
\section{Related Work}
\label{sec:related}
%--------------------------------------------------------------------

\textbf{X.509 and PKI.}  RFC~5280~\cite{rfc5280} places a CA hierarchy
beneath individual keys.  CAs can mis-issue certificates, be compromised, or
be coerced, violating the requirement that no institutional key sits beneath
the individual~\cite{goldreich2001}.  Identity Bedrock removes this layer
entirely.

\textbf{Decentralised Identifiers (DIDs).}  The W3C DID specification
\cite{did2022} defines a generic framework for self-sovereign identifiers.
Identity Bedrock is compatible with the DID model at the identifier layer but
provides a stronger formal security foundation and a concrete economic
finality argument absent from the DID core specification.

\textbf{FIDO2 / WebAuthn.}  The FIDO2 standard~\cite{webauthn2021} eliminates
passwords using device-bound asymmetric keys, closely aligning with the
identity primitive of Definition~1.  However, FIDO2 relies on attestation
certificates from hardware vendors, reintroducing an institutional trust
anchor.  Identity Bedrock requires no attestation layer.

\textbf{Sign-In with Ethereum (EIP-4361).}  EIP-4361~\cite{eip4361} binds
identity to Ethereum account keys and uses Ethereum state as the directory.
This couples identity finality to Ethereum consensus and governance, and
restricts the signature scheme to secp256k1.  Identity Bedrock generalises
this pattern to arbitrary EUF-CMA schemes and arbitrary append-only
directories.

\textbf{ENS and Handshake.}  The Ethereum Name Service (ENS) and
Handshake~\cite{handshake} provide decentralised naming but do not define a
formal authentication substrate.  They can be viewed as concrete
instantiations of the directory layer $\mathcal{D}$, compatible with the
architecture presented here.

\textbf{Verifiable Credentials.}  The W3C Verifiable Credentials (VC)
model~\cite{vc2019} supports issuer-signed claims, which introduce issuer
keys as trust anchors.  Identity Bedrock operates one layer below VCs: it
provides the identity binding from which a VC holder's key can be derived.

%--------------------------------------------------------------------
\section{Discussion}
\label{sec:discussion}
%--------------------------------------------------------------------

\textbf{What the bedrock provides.}  A verifiable, unforgeable binding between
an identifier and a public key, secured by standard cryptographic assumptions
and a quantified economic finality argument.  All higher-layer properties
described in Section~\ref{sec:derived} inherit this security basis.

\textbf{What is out of scope.}  The model does not address:
\begin{itemize}[leftmargin=*]
  \item \emph{Endpoint compromise and malware.}  An adversary with access to
    the device holding $sk$ trivially breaks authentication regardless of
    the protocol.
  \item \emph{Social engineering and phishing.}  Out-of-band attacks that
    induce the legitimate key holder to sign adversarial challenges.
  \item \emph{Legal and regulatory coercion.}  Compelled key disclosure is
    orthogonal to cryptographic security.
\end{itemize}
These are operational and social concerns that require additional layers and
are deliberately out of scope.

\textbf{Comparison with zero-knowledge approaches.}  The authentication
protocol in Section~4.1 reveals the challenge--response pair $(c, \sigma)$,
which may allow linkability across sessions if $pk$ is public.  Anonymous
credential schemes~\cite{goldreich2001} can provide unlinkability at the cost
of additional complexity; this is a natural extension of the bedrock rather
than a contradiction of it.

\textbf{Post-quantum migration.}  Ed25519 is vulnerable to Shor's algorithm.
The intended migration path is ML-DSA-65 (CRYSTALS-Dilithium3,
FIPS~204~\cite{fips204}):
\begin{itemize}[leftmargin=*]
  \item Replace Ed25519 with ML-DSA-65 (public key 1952~bytes; accommodated
    by the EmerNVS value field).
  \item The WORM \texttt{<keys>} schema~\cite{worm} natively supports multiple
    algorithms; migration adds \texttt{<key algo="mldsa65">} without
    modifying existing records.
  \item Directory anchoring, NESS Fiber structure, and the authentication
    protocol are unchanged.
  \item OpenSSL~3.5 (April 2025) includes native ML-DSA support in mainline
    Armbian Linux.
\end{itemize}
The directory and connectivity layers are post-quantum agnostic; only the
signature scheme at Layer~1 requires migration.

\textbf{Key rotation.}  The directory model supports key rotation by
registering a new binding $(id, pk')$ and revoking the old one through a
signed revocation record.  The append-only property means revocations are
themselves permanent and publicly verifiable.

%--------------------------------------------------------------------
\section{Conclusion}
%--------------------------------------------------------------------

Identity Bedrock demonstrates that a minimal cryptographic substrate---an
EUF-CMA signature scheme and an append-only directory with quantified
finality---is sufficient not only for authentication but for enabling a full
suite of system capabilities without introducing institutional trust anchors.

The formal model, security proof, and economic attack-cost analysis establish
that the only meaningful attacks on the system require either breaking standard
hardness assumptions or paying a publicly quantifiable economic cost.  The
system has been deployed and operational since 2020 on commodity hardware,
validating the practical feasibility of the construction.

By deriving naming, content addressing, deterministic execution, connectivity,
and economic interaction from a single identity primitive, the architecture
reduces conceptual complexity while preserving compatibility with existing
protocols and standards.

%--------------------------------------------------------------------
\section*{Acknowledgments}
%--------------------------------------------------------------------

The author thanks Oleg Khovayko for foundational Emercoin infrastructure
work and the hack51 attack-cost framework~\cite{hack51}; Aleksej Sokolov
(Alex014) for designing the WORM format~\cite{worm} and implementing the
PrivateNess toolchain (privatenesstools, pyuheprng) from the author's
specifications; Brandon ``Synth'' Smietana and Moses Narrow
(\texttt{moe\_narrow} / \texttt{0pcom}), Skywire Linux packaging lead, for
the Skywire mesh transport and Fiber consensus infrastructure on which NESS
is built; and the PrivateNess Network community, whose collective work since
the project's public founding on Bitcointalk on September~1,
2020~\cite{bitcointalk} has made this deployed system possible.  The author's
own engagement with Emercoin infrastructure spans nine years and provides the
operational foundation on which Identity Bedrock is built.

%--------------------------------------------------------------------
\begin{thebibliography}{99}

\bibitem{rompel1990}
  J.~Rompel,
  ``One-way functions are necessary and sufficient for secure signatures,''
  in \textit{Proc.\ 22nd Annual ACM STOC}, pp.~387--394, 1990.

\bibitem{haitner2010}
  I.~Haitner, T.~Holenstein, O.~Reingold, S.~Vadhan, and H.~Wee,
  ``Universal one-way hash functions via inaccessible entropy,''
  in H.~Gilbert (Ed.), \textit{EUROCRYPT 2010},
  LNCS vol.~6110, pp.~616--637, Springer, 2010.

\bibitem{naoryung1989}
  M.~Naor and M.~Yung,
  ``Universal one-way hash functions and their cryptographic applications,''
  in \textit{Proc.\ 21st Annual ACM STOC}, pp.~33--43, 1989.

\bibitem{katz2005}
  J.~Katz and C.-Y.~Koo,
  ``On constructing universal one-way hash functions from arbitrary one-way
  functions,''
  Cryptology ePrint Archive, Report 2005/328, 2005.
  \url{https://eprint.iacr.org/2005/328}

\bibitem{goldreich2001}
  O.~Goldreich,
  \textit{Foundations of Cryptography, Volume~1: Basic Tools},
  Cambridge University Press, 2001.

\bibitem{eddsa}
  S.~Josefsson and I.~Liusvaara,
  ``Edwards-Curve Digital Signature Algorithm (EdDSA),''
  RFC~8032, IETF, January 2017.
  \url{https://www.rfc-editor.org/rfc/rfc8032}

\bibitem{fips204}
  National Institute of Standards and Technology,
  ``Module-Lattice-Based Digital Signature Standard (ML-DSA),''
  FIPS~204, August 2024.
  \url{https://doi.org/10.6028/NIST.FIPS.204}

\bibitem{uheprng}
  S.~Gibson,
  ``Ultra-High Entropy Pseudo-Random Number Generator (UHEPRNG),''
  Gibson Research Corporation.
  \url{https://www.grc.com/otg/uheprng.htm}

\bibitem{emercoin}
  Emercoin Development Team,
  ``Emercoin: hybrid blockchain with merged mining,''
  \url{https://emercoin.com}, 2013--2026.

\bibitem{hack51}
  O.~Khovayko,
  ``Estimating the cost of a 51\% attack on Emercoin: the hack51 coefficient''
  [In Russian],
  Dzen / Yandex, 2021.
  \url{https://dzen.ru/a/YUotZkN9nAeElPEQ}

\bibitem{nuzzi2024}
  L.~Nuzzi et al.,
  ``Breaking BFT: quantifying the cost to attack Bitcoin and Ethereum,''
  Coin Metrics, 2024.

\bibitem{skycoin}
  B.~Smietana (``Synth'') et al.,
  ``Skycoin: Skywire and Fiber consensus,''
  Skycoin Development Team, 2014--2026.
  \url{https://www.skycoin.com}
  Linux packaging: M.~Narrow (\texttt{moe\_narrow} / \texttt{0pcom}).

\bibitem{worm}
  A.~Sokolov (Alex014) / NESS-Network,
  ``WORM --- World Object MaRkup language,''
  format designed by A.~Sokolov; deployed within the PrivateNess ecosystem.
  \url{https://github.com/ness-network/worm}

\bibitem{onchain}
  PrivateNess Network,
  ``PrivateNess.Network --- canonical on-chain identity record,''
  Emercoin NVS DPO, anchored in block~\#445669,
  November~1, 2020, 22:16:45~UTC.  Operational 2020--2026.
  \url{https://explorer.emercoin.com/dpo/PrivateNESS.Network}

\bibitem{bitcointalk}
  PrivateNess Network,
  ``PrivateNess Network (NESS) Privacy CX Fiber Coin,''
  Bitcointalk Announcements (Altcoins), September~1, 2020.
  \url{https://bitcointalk.org/index.php?topic=5272741.0}

\bibitem{rfc5280}
  D.~Cooper et al.,
  ``Internet X.509 Public Key Infrastructure Certificate and CRL Profile,''
  IETF RFC~5280, May 2008.

\bibitem{did2022}
  M.~Sporny et al.,
  ``Decentralized Identifiers (DIDs) v1.0,''
  W3C Recommendation, July 2022.
  \url{https://www.w3.org/TR/did-core/}

\bibitem{webauthn2021}
  D.~Balfanz et al.,
  ``Web Authentication: An API for Accessing Public Key Credentials Level~2,''
  W3C Recommendation, April 2021.
  \url{https://www.w3.org/TR/webauthn-2/}

\bibitem{eip4361}
  B.~Talsma et al.,
  ``EIP-4361: Sign-In with Ethereum,''
  Ethereum Improvement Proposal, October 2021.
  \url{https://eips.ethereum.org/EIPS/eip-4361}

\bibitem{vc2019}
  M.~Sporny, D.~Longley, and D.~Chadwick,
  ``Verifiable Credentials Data Model 1.0,''
  W3C Recommendation, November 2019.
  \url{https://www.w3.org/TR/vc-data-model/}

\bibitem{handshake}
  Handshake Authors,
  ``Handshake: Decentralized Certificate Authority and Naming,''
  Technical Paper, 2018.
  \url{https://handshake.org/files/handshake.txt}

\bibitem{ipfs}
  J.~Benet,
  ``IPFS---Content Addressed, Versioned, P2P File System,''
  arXiv:1407.3561, 2014.

\end{thebibliography}

\end{document}
